\documentclass[11pt,a4paper]{article}
\usepackage[utf8]{inputenc}
\usepackage[T1]{fontenc}
\usepackage{graphicx}
\usepackage{xcolor}
\usepackage{hyperref}
\usepackage{geometry}
\usepackage{fancyhdr}
\usepackage{enumitem}
\setlength{\parskip}{6pt}
\setlength{\parindent}{0pt}

\geometry{margin=2.2cm}
\hypersetup{colorlinks=true,linkcolor=black,urlcolor=blue!50!black}

\definecolor{nexagreen}{HTML}{00D4AA}
\definecolor{nexadark}{HTML}{0B0B0B}
\definecolor{nexagray}{HTML}{666666}

\pagestyle{fancy}
\fancyhf{}
\fancyhead[L]{\small\color{nexagray} NexaPay --- Whitepaper Technique}
\fancyhead[R]{\small\color{nexagray} La Banque de Tunisie}
\fancyfoot[C]{\thepage}
\renewcommand{\headrulewidth}{0.4pt}
\renewcommand{\familydefault}{\sfdefault}

\begin{document}

% ─── COVER ───
\begin{center}
\vspace*{3cm}
{\Huge\bfseries NexaPay}\\[0.3cm]
{\Large\color{nexagreen}\bfseries Plateforme de Paiement Num\'erique}\\[0.6cm]
{\large Whitepaper Technique}\\[0.3cm]
{\large\color{nexagray} Pr\'epar\'e pour}\\[0.2cm]
{\Large\bfseries La Banque de Tunisie}\\[1cm]
{\color{nexagreen}\rule{10cm}{2pt}}\\[0.8cm]
{\large\color{nexagray} Mai 2026}\\[1.5cm]
{\normalsize\bfseries Glitch Inc --- BackendGlitch Division}\\[0.2cm]
{\small\color{nexagray} backendglitch.com $|$ contact@backendglitch.com}\\[0.2cm]
{\small\color{nexagray} nexapay.space $|$ sandbox.nexapay.space}
\end{center}
\newpage

% ─── TOC ───
\tableofcontents
\newpage

% ═══════════════════════════════════════════════════════
\section{Pr\'esentation de NexaPay}

NexaPay est une plateforme fintech compl\`ete con\c{c}ue pour le march\'e tunisien. Elle combine un portefeuille num\'erique, une passerelle de paiement pour commer\c{c}ants, et un moteur blockchain propri\'etaire avec consensus BFT \`a 3 validateurs.

\subsection{URLs Publiques}

\begin{description}[leftmargin=*,style=nextline]
    \item[\href{https://nexapay.space}{nexapay.space}] Page d'accueil publique
    \item[\href{https://sandbox.nexapay.space}{sandbox.nexapay.space}] Plateforme principale (dashboard, envois, carte)
    \item[\href{https://auth.nexapay.space}{auth.nexapay.space}] Authentification (login, inscription)
    \item[\href{https://backend.nexapay.space}{backend.nexapay.space}] API Backend
    \item[\href{https://nexapay.space/docs}{nexapay.space/docs}] Documentation API \& SDK
\end{description}

\subsection{Architecture}

\begin{itemize}[leftmargin=*]
    \item \textbf{Frontend} --- Next.js 16 avec routage par sous-domaines
    \item \textbf{Backend} --- Rust/Axum avec JWT, cookies httpOnly, AES-256-GCM
    \item \textbf{Blockchain} --- Cha\^ine append-only, 3 validateurs BFT, signatures Ed25519
    \item \textbf{Base de donn\'ees} --- PostgreSQL 16 + SQLite + Sled
    \item \textbf{SDK} --- \texttt{@nexapay/node-sdk} sur npm
\end{itemize}

\newpage

% ═══════════════════════════════════════════════════════
\section{S\'ecurit\'e et Protection des Donn\'ees}

La s\'ecurit\'e est au c\oe{}ur de l'architecture NexaPay. Chaque composant a \'et\'e con\c{c}u avec les meilleures pratiques de l'industrie financi\`ere.

\subsection{Protection des Donn\'ees Utilisateur}

\textbf{Donn\'ees personnelles :}
\begin{itemize}[leftmargin=*]
    \item Les donn\'ees sensibles (CIN, num\'ero de t\'el\'ephone, email) sont stock\'ees dans PostgreSQL avec acc\`es strictement contr\^ol\'e
    \item Les recherches de comptes ne retournent que les 4 derniers chiffres du t\'el\'ephone et du CIN --- jamais les donn\'ees compl\`etes
    \item Toutes les communications entre le navigateur et le serveur passent par HTTPS/TLS 1.3
\end{itemize}

\textbf{Donn\'ees bancaires :}
\begin{itemize}[leftmargin=*]
    \item Les num\'eros de carte Visa et CVV sont chiffr\'es avec AES-256-GCM (standard militaire)
    \item Seuls les 4 derniers chiffres de la carte sont visibles --- le num\'ero complet n'est jamais affich\'e
    \item Les cl\'es priv\'ees Ed25519 des utilisateurs sont chiffr\'ees avec une cl\'e d\'eriv\'ee de leur PIN personnel via Argon2id
    \item Chaque utilisateur poss\`ede un \textbf{salt al\'eatoire unique} pour le hachage de son PIN --- rendant les attaques par dictionnaire impossibles
\end{itemize}

\textbf{Sessions et Authentification :}
\begin{itemize}[leftmargin=*]
    \item Double authentification : PIN (6 chiffres) + code OTP envoy\'e par SMS
    \item Les sessions utilisent des cookies \textbf{httpOnly} (inaccessibles au JavaScript, immunis\'es contre le vol par XSS)
    \item Les tokens JWT utilisateur et administrateur sont sign\'es avec des \textbf{cl\'es cryptographiques distinctes}
    \item Les sessions peuvent \^etre r\'evoques c\^ot\'e serveur \`a tout moment
    \item Protection contre les attaques par force brute : rate limiting par adresse IP sur toutes les routes d'authentification
\end{itemize}

\textbf{Transactions et Blockchain :}
\begin{itemize}[leftmargin=*]
    \item Chaque transaction est enregistr\'ee dans un registre blockchain immuable (append-only)
    \item Les signatures Ed25519 garantissent la non-r\'epudiation --- aucune transaction ne peut \^etre falsifi\'ee
    \item Le consensus \`a 3 validateurs assure qu'aucune entit\'e seule ne peut modifier l'\'etat de la cha\^ine
    \item Les webhooks (notifications commer\c{c}ants) sont sign\'es avec HMAC-SHA256 pour garantir leur authenticit\'e
\end{itemize}

\subsection{Pourquoi 3 Validateurs (BFT) ?}

NexaPay utilise un consensus \textbf{Byzantine Fault Tolerant} \`a 3 validateurs. Contrairement aux syst\`emes centralis\'es (un seul serveur) ou aux blockchains publiques (lentes et co\^uteuses) :

\begin{itemize}[leftmargin=*]
    \item \textbf{Aucun point de d\'efaillance unique} --- Si un validateur tombe en panne, les deux autres continuent de fonctionner
    \item \textbf{Validation crois\'ee} --- Chaque bloc doit \^etre sign\'e par les 3 validateurs avant d'\^etre ajout\'e \`a la cha\^ine
    \item \textbf{Transactions instantan\'ees} --- Pas de minage, pas de proof-of-work. Les transferts sont appliqu\'es en moins d'une seconde
    \item \textbf{Co\^ut z\'ero} --- Pas de frais de gas, pas de commission blockchain. Les transferts entre utilisateurs sont gratuits
    \item \textbf{Souverainet\'e} --- L'infrastructure est h\'eberg\'ee sur nos propres serveurs, pas sur un cloud public \'etranger
\end{itemize}

\subsection{Flux d'Authentification}

\begin{enumerate}[leftmargin=*]
    \item Utilisateur entre t\'el\'ephone + PIN (6 chiffres)
    \item V\'erification PIN avec Argon2id (sel al\'eatoire unique)
    \item Envoi OTP 6 chiffres par SMS (Twilio)
    \item V\'erification OTP en temps constant (r\'esistant aux attaques temporelles)
    \item Session \'etablie avec cookie httpOnly + JWT
\end{enumerate}

\newpage

% ═══════════════════════════════════════════════════════
\section{Processus KYC (Know Your Customer)}

\textbf{Note :} Le module KYC complet est disponible mais \textbf{d\'esactiv\'e pour la d\'emo}. Les utilisateurs sont auto-v\'erifi\'es en mode sandbox. Le flux ci-dessous sera activ\'e en production.

\subsection{\'Etapes du KYC}

\begin{enumerate}[leftmargin=*]
    \item \textbf{Soumission des documents}
    \begin{itemize}
        \item Photographie du \textbf{recto} et \textbf{verso} de la CIN
        \item Upload s\'ecuris\'e avec validation (PNG, JPG, max 5 Mo)
        \item Stockage chiffr\'e des documents
    \end{itemize}

    \item \textbf{Extraction automatique (OCR)}
    \begin{itemize}
        \item Nom, pr\'enom, num\'ero CIN, date de naissance, date d'\'emission, lieu de naissance
        \item Croisement avec le formulaire d'inscription
    \end{itemize}

    \item \textbf{Test de vivacit\'e (Liveness Detection)}
    \begin{itemize}
        \item V\'erification via cam\'era (clignement, rotation t\^ete, sourire)
        \item Score de correspondance faciale entre photo CIN et selfie
        \item Protection anti-photo et anti-vid\'eo
    \end{itemize}

    \item \textbf{V\'erification et approbation}
    \begin{itemize}
        \item Score de vivacit\'e $>$ seuil requis ET donn\'ees CIN correspondantes
        \item Compte marqu\'e \texttt{kyc\_status = verified}
        \item Acc\`es complet \`a la plateforme
    \end{itemize}
\end{enumerate}

\subsection{Donn\'ees KYC Stock\'ees}

\begin{description}[leftmargin=*]
    \item[\texttt{kyc\_cin\_front\_path}] Chemin recto CIN
    \item[\texttt{kyc\_cin\_back\_path}] Chemin verso CIN
    \item[\texttt{kyc\_face\_match\_score}] Score correspondance faciale
    \item[\texttt{kyc\_liveness\_passed}] R\'esultat test vivacit\'e
    \item[\texttt{kyc\_verified\_at}] Date/heure v\'erification
\end{description}

\subsection{Conformit\'e R\'eglementaire}

Le processus est conforme avec la r\'eglementation BCT, les directives LCB-FT, et la Loi n$^\circ$ 2004-63.

\newpage

% ═══════════════════════════════════════════════════════
\section{Parcours Utilisateur}

\subsection{Inscription (6 \'etapes)}

\begin{enumerate}[leftmargin=*]
    \item \textbf{T\'el\'ephone} --- Num\'ero tunisien (+216 XX XXX XXX)
    \item \textbf{Informations personnelles} --- Nom, pr\'enom, email, date de naissance
    \item \textbf{Identit\'e} --- Num\'ero CIN et date d'\'emission
    \item \textbf{Contrat} --- Signature \'electronique du contrat (Loi n$^\circ$ 2002-50)
    \item \textbf{PIN} --- Code 6 chiffres pour transactions
    \item \textbf{Activation} --- IBAN/RIB tunisien + carte Visa virtuelle + 150 TND offerts
\end{enumerate}

\subsection{Fonctionnalit\'es Utilisateur}

\begin{itemize}[leftmargin=*]
    \item \textbf{Portefeuille} --- Solde TND, IBAN/RIB, carte Visa virtuelle
    \item \textbf{Envoi d'argent} --- Transferts instantan\'es et gratuits entre utilisateurs
    \item \textbf{Virement bancaire} --- Vers n'importe quel RIB tunisien
    \item \textbf{Rechargement} --- D\'ep\^ot esp\`eces en agence bancaire partenaire
    \item \textbf{Paiement en ligne} --- Checkout NexaPay chez les commer\c{c}ants
    \item \textbf{Historique} --- Transactions on-chain, v\'erifiables
    \item \textbf{Carte virtuelle} --- Achats en ligne, gel/d\'egel
\end{itemize}

\subsection{D\'ep\^ot d'Argent en Agence}

\begin{enumerate}[leftmargin=*]
    \item Le client se rend dans une agence de \textbf{La Banque de Tunisie}
    \item Il fournit son num\'ero de t\'el\'ephone ou son adresse NexaPay
    \item L'agent bancaire v\'erifie l'identit\'e du client
    \item Le d\'ep\^ot en esp\`eces est effectu\'e au guichet
    \item L'agent cr\'edite le compte via l'interface administrateur
    \item Le solde est mis \`a jour en temps r\'eel
\end{enumerate}

\newpage

% ═══════════════════════════════════════════════════════
\section{Programme Agent / Commer\c{c}ant}

\subsection{Devenir Agent}

\begin{enumerate}[leftmargin=*]
    \item \textbf{Candidature} --- Formulaire avec nom commercial, type d'activit\'e, matricule fiscal, adresse, volume estim\'e. Upload licence commerciale (PDF) et extrait RNE (PDF)
    \item \textbf{Scoring} --- \'Evaluation automatique du risque
    \item \textbf{Approbation} --- Automatique en sandbox, r\'evision manuelle en production
    \item \textbf{Activation} --- Cl\'e API unique, dashboard agent, limite de volume
\end{enumerate}

\subsection{Capacit\'es de l'Agent}

\begin{itemize}[leftmargin=*]
    \item \textbf{Liens de paiement} --- Cr\'eation de liens/QR codes pour accepter des paiements
    \item \textbf{API REST} --- Int\'egration via API ou SDK Node.js
    \item \textbf{Remboursements} --- Partiels ou totaux
    \item \textbf{Solde} --- Disponible, volume brut, retraits
    \item \textbf{Retraits} --- Transfert solde commercial $\rightarrow$ portefeuille personnel (on-chain)
    \item \textbf{Webhooks} --- Notifications temps r\'eel (succ\`es/\'echec)
    \item \textbf{Statistiques} --- Volume total, taux de succ\`es, transactions du jour
    \item \textbf{Export} --- CSV des transactions et clients
\end{itemize}

\subsection{Int\'egration Technique}

Trois modes d'int\'egration :

\begin{enumerate}[leftmargin=*]
    \item \textbf{No-Code} --- Liens/QR codes depuis le dashboard
    \item \textbf{SDK Node.js} --- \texttt{npm install @nexapay/node-sdk}
    \item \textbf{REST API} --- Endpoints HTTP avec cl\'e API
\end{enumerate}

\subsubsection{Exemple API}

\begin{verbatim}
curl -X POST https://backend.nexapay.space/gateway/v1/intents \
  -H "Content-Type: application/json" \
  -H "X-API-Key: nxp_developer_CLE_API" \
  -d '{"amount":42000,"currency":"TND","description":"Cmd #42",
       "success_webhook_url":"https://maboutique.tn/hook/success"}'
\end{verbatim}

\newpage

% ═══════════════════════════════════════════════════════
\section{Blockchain et Tra\c{c}abilit\'e}

\begin{itemize}[leftmargin=*]
    \item \textbf{Consensus BFT} --- 3 validateurs, quorum 3/3, signatures Ed25519
    \item \textbf{Registre immuable} --- Toutes les transactions on-chain (append-only)
    \item \textbf{12 types de transactions} --- Transfer, AccountCreate, EsignAccount, InvoiceAnchor, etc.
    \item \textbf{Ancrage} --- Contrats sign\'es et factures hash\'es et ancr\'es on-chain
    \item \textbf{Mempool} --- 500 transactions max, expiration 5 minutes
    \item \textbf{Persistance} --- Sled (blocs) + SQLite (snapshots)
\end{itemize}

\subsection{Propri\'et\'es}

\begin{itemize}[leftmargin=*]
    \item \textbf{Int\'egrit\'e} --- Chaque bloc r\'ef\'erence le hash du bloc pr\'ec\'edent
    \item \textbf{Non-r\'epudiation} --- Signatures Ed25519 sur transactions utilisateur
    \item \textbf{Auditabilit\'e} --- Toute transaction v\'erifiable ind\'ependamment
    \item \textbf{R\'esilience} --- 3 validateurs redondants
\end{itemize}

\newpage

% ═══════════════════════════════════════════════════════
\section{Pourquoi NexaPay ? --- Avantages Concurrentiels}

\subsection{Comparaison avec les E-Wallets Existants}

\begin{center}
\begin{tabular}{|p{3.8cm}|p{4.2cm}|p{4.2cm}|}
\hline
\textbf{} & \textbf{E-Wallets Traditionnels} & \textbf{NexaPay} \\
\hline
\textbf{Infrastructure} & Serveur centralis\'e unique & 3 validateurs BFT redondants \\
\hline
\textbf{Vitesse transferts} & 1--3 jours ouvr\'es (virement) & Instantan\'e ($<$1 seconde) \\
\hline
\textbf{Frais transferts} & 0.5--2\% par transaction & Gratuit entre utilisateurs \\
\hline
\textbf{Tra\c{c}abilit\'e} & Base de donn\'ees interne & Blockchain immuable v\'erifiable \\
\hline
\textbf{IBAN/RIB} & D\'epend de la banque partenaire & G\'en\'er\'e automatiquement \\
\hline
\textbf{Carte virtuelle} & Parfois disponible & Visa virtuelle incluse \\
\hline
\textbf{API Commer\c{c}ants} & Limit\'ee ou absente & REST API + SDK npm \\
\hline
\textbf{Webhooks} & Rare & Sign\'es HMAC-SHA256 avec retry \\
\hline
\textbf{KYC} & Manuel (agence) & Automatis\'e (OCR + liveness) \\
\hline
\textbf{Contrat \'electronique} & Papier & Sign\'e + ancr\'e blockchain (Loi 2002-50) \\
\hline
\textbf{Chiffrement} & Variable & AES-256-GCM bout en bout \\
\hline
\end{tabular}
\end{center}

\subsection{Ce Qui Rend NexaPay Unique}

\begin{enumerate}[leftmargin=*]
    \item \textbf{Blockchain priv\'ee \`a 3 validateurs} --- Ni centralis\'e (risque de panne unique), ni blockchain publique (lente et ch\`ere). Un compromis id\'eal pour la finance : rapidit\'e, s\'ecurit\'e, et co\^ut z\'ero.

    \item \textbf{IBAN/RIB tunisien g\'en\'er\'e automatiquement} --- Chaque compte NexaPay re\c{c}oit un IBAN et RIB tunisien. Les virements bancaires standards fonctionnent sans adaptation.

    \item \textbf{Double r\^ole utilisateur/commer\c{c}ant} --- Un m\^eme compte peut \^etre portefeuille personnel ET passerelle de paiement. Pas besoin de deux applications distinctes.

    \item \textbf{Ancrage blockchain des documents} --- Les contrats d'ouverture de compte et les factures sont hash\'es et enregistr\'es on-chain. Toute autorit\'e peut v\'erifier l'authenticit\'e d'un document de mani\`ere ind\'ependante.

    \item \textbf{Con\c{c}u pour la Tunisie} --- Conforme \`a la r\'eglementation BCT, formats tunisiens (CIN, RIB 20 chiffres, num\'eros +216), interface en fran\c{c}ais, support des d\'ep\^ots en esp\`eces en agence.

    \item \textbf{Int\'egration simple pour les commer\c{c}ants} --- 3 modes : liens de paiement (no-code), SDK npm, ou REST API. Le commer\c{c}ant choisit selon ses capacit\'es techniques.
\end{enumerate}

\newpage

% ═══════════════════════════════════════════════════════
\section{Contact}

\begin{center}
{\Large\bfseries Glitch Inc --- BackendGlitch Division}

\vspace{0.5cm}
\begin{tabular}{ll}
\textbf{Site web} & \href{https://backendglitch.com}{backendglitch.com} \\
\textbf{Email} & \href{mailto:contact@backendglitch.com}{contact@backendglitch.com} \\
\textbf{GitHub} & \href{https://github.com/Samer-Gassouma/NexaPay}{github.com/Samer-Gassouma/NexaPay} \\
\textbf{SDK npm} & \href{https://www.npmjs.com/package/@nexapay/node-sdk}{npmjs.com/package/@nexapay/node-sdk} \\
\end{tabular}

\vspace{0.8cm}
\begin{tabular}{ll}
\textbf{Landing} & \href{https://nexapay.space}{nexapay.space} \\
\textbf{Sandbox} & \href{https://sandbox.nexapay.space}{sandbox.nexapay.space} \\
\textbf{API Backend} & \href{https://backend.nexapay.space}{backend.nexapay.space} \\
\textbf{Documentation} & \href{https://nexapay.space/docs}{nexapay.space/docs} \\
\end{tabular}

\vspace{1.5cm}
{\color{nexagreen}\rule{8cm}{1pt}}

\vspace{0.5cm}
{\large\bfseries\color{nexagreen} Merci de votre attention.}

\vspace{0.3cm}
{\normalsize\color{nexagray} Disponible pour d\'emonstration en direct et questions techniques.}

\vspace{0.5cm}
{\small\color{nexagray} Glitch Inc --- BackendGlitch Division \copyright\ 2026}
\end{center}

\end{document}
